% 20260529_Assingment_Probability_Statistics
\documentclass[dvipdfmx]{jsreport}
%preamble
\usepackage{soupreamble}

\title{確率及び統計学特論A 第3回課題}
\author{M1 6012 川村聡一郎}
\date{}


\begin{document}
\maketitle

\begin{tcolorbox}
	Youngの不等式を示せ．\\
	$p, q> 1, \frac{1}{p}+ \frac{1}{q}= 1$とする．このとき$a, b\ge 0$に対して，
\begin{equation*}
	ab\le \frac{1}{p}a^p+ \frac{1}{q}b^q
\end{equation*}
\end{tcolorbox}

\begin{souanswer} %答案はここに書く
	$f(x)= -\log{x}$とする．$f''(x)= \frac{1}{x^2}> 0$なので$(0, \infty)$で凸．凸関数の定義より，
\begin{align*}
	-\log{\left(\frac{1}{p}a^p+ \frac{1}{q}b^q\right)}
	&\le -\frac{1}{p}\log{a^p}- \frac{1}{q}\log{b^q}\\
	&= -\log{ab}
\end{align*}
	すなわち$\log{\left(\frac{1}{p}a^p+ \frac{1}{q}b^q\right)}\le \log{ab}$．\\
	底は正だから大小関係は変わらず，Youngの不等式を得る．
\end{souanswer}
\end{document}